\scanpage{106}
% Approximately lines 1--2; continuation from scan 105.
\(u'_{\Gamma}=\operatorname{int}_{\mathcal E}(g)\circ u_{\Gamma}\), and we finish as
above, with uniqueness of \(g\).

% Approximately lines 3--8.
Now take an \inserted{\uline{injective}} outer homomorphism
\[
\deleted{\unsure{\lambda}}\Lambda=[\kappa]:T\longrightarrow\pi
\qquad (T\text{ commutative}),
\]
% Written directly under the letter Λ.
\margin{(initial of ``loops'')}
% Oblique note in the upper-left margin, tied by a curve to "injectif".
\margin{Injectivity in the case we are interested in follows from the anabelian hypothesis.}
% "supposons" stands at the start of the line, with an arrow to "que".
suppose that (for \(\kappa\in[\kappa]\))
\[
\operatorname{Centr}_\pi\kappa=\kappa(T)
\]
(which does not depend \unsure{on the|on the} choice of \(\kappa\) in
\([\kappa]\)). I will then define a
% Approximately lines 9--17.
\deleted{category} connected abelian groupoid \inserted{\(\Pi_\Lambda\)}, and an
\unsure{isomorphism} of its \(\pi_{\unsure{1}}\) with \(T\) \deleted{An object}
(hence a topos, which plays the role of
\(\widehat{\Pi}_{\unsure{D|D^*}}\)). An object will be a \(\kappa\in[\kappa]\).
A homomorphism from \(\kappa\) to \(\kappa'\) will be an
element \(g\in\operatorname{Transp}_\pi(\kappa,\kappa')\). \deleted{Let \(g\)}
The composition of homomorphisms obvious. We
do indeed get a connected groupoid, in
which \(\operatorname{Aut}(\kappa)=\operatorname{Centr}_\pi(\kappa)\overset{\inserted{\text{by hypothesis}}}{=}\kappa(T)\unsure{\xleftarrow[\sim]{\kappa}}T\).

% Approximately lines 18--21; continuation on scan 107.
% The words from "sur les topos" to "B_Λ → B_π" are squeezed between two lines.
OK. Of course, we have a homomorphism
\(\Pi_\Lambda\longrightarrow\deleted{\unsure{\pi}}\) one-point groupoid, hence
\unsure{on} the \unsure{classifying|classifying} topoi \inserted{defined by \(\pi\)} \inserted{\(B_\Lambda\longrightarrow B_\pi\)}.
We have to see the behavior of this
% Lower-left additional handwritten pagination; hand/status uncertain.
\margin{105}

\scanpage{107}
\mspage{52}
% Approximately lines 1--7; continuation from scan 106.
construction under outer homomorphisms. So let
\[
\Theta=f_\pi\deleted{\unsure{\Theta}}=[u]:\pi'\longrightarrow\pi
\]
be an outer homomorphism such that \(\operatorname{Centr}_\pi(\Theta)=1\),
and \(\Lambda':T\longrightarrow\pi'\) an injective outer homomorphism
such that \(\kappa'\in\Lambda'\Rightarrow\operatorname{Centr}_{\pi'}(\kappa')=\operatorname{Im}\kappa'\)
and let us consider the composite
\[
\deleted{\unsure{\Lambda=}}f_\pi\circ\Lambda':T\longrightarrow\pi.
\]
% Approximately lines 8--12.
Let \(d\in\mathbb Z\) and consider
\[
\Lambda=f_\pi\circ\Lambda'\circ(d.\operatorname{id}_T)
\]
\inserted{suppose that \(\kappa\in\Lambda\Rightarrow\operatorname{Centr}_\pi\kappa=\operatorname{Im}\kappa\)}
We are going to define a homomorphism from \(\Pi_{\Lambda'}\) to \(\Pi_\Lambda\),
hence a homomorphism of Topoi \(B_{\Lambda'}\longrightarrow B_\Lambda\),
and a commutation datum \inserted{\(\alpha\)} for the
\diagram{\begin{tikzcd}[ampersand replacement=\&,column sep=4em,row sep=3.8em]
\& B_{\Lambda'} \arrow[r,"f_D"] \arrow[d] \& B_\Lambda \arrow[d] \\
(*) \& \unsure{(B,\pi')} \arrow[r,"{\begin{gathered}\deleted{\unsure{B(f_\pi)}}\\[-2pt]\inserted{B(f_\pi)\ \text{defined above}}\end{gathered}}"']
\arrow[ur,bend left=32,shorten <=2.1em,shorten >=2.1em,"\alpha"']
\& \unsure{(B,\pi)}
\end{tikzcd}}
\margin{\illegible{2} \unsure{obviously} \illegible{5} \deleted{\unsure{of Topoi}} \illegible{2}.}
\deleted{\illegible{2} \unsure{to the} hypotheses \illegible{1} on \(\Lambda',\Theta,d\)}

% Approximately lines 13--19; continuation on scan 108.
\deleted{Let} \(\kappa':T\longrightarrow\pi'\) an object of \(\Pi_{\Lambda'}\),
we want to define (up to unique isomorphism) an
object \(\kappa\) of \(\Pi_\Lambda\). For every \(u\in f_\pi=\Theta\),
% The κ in the next formula carries no prime in the ink.
we consider \(u\circ\kappa\circ(d.\operatorname{id}_T)\), there is between them
a transitive system of isomorphisms for \(u\) varying
in \(\Theta\), we can identify them with one another.
% Lower-left additional handwritten pagination; hand/status uncertain.
\margin{106}

\scanpage{108}
% Approximately lines 1--7.
The functoriality of this object with respect to
varying \(\kappa'\) is obvious: [we can say that
\(u\in\Theta\) defines a homomorphism of groupoids
\[
u_\Lambda:\Pi_\Lambda\longrightarrow\Pi_{\Lambda'},
\]
and between these
there is a transitive system of isomorphisms \unsure{\}}. In fact,
for fixed \(u\) we have a \deleted{system} commutative
diagram of homomorphisms of groupoids
\diagram{\begin{tikzcd}[ampersand replacement=\&]
\Pi_{\Lambda'} \arrow[r,"u_\Lambda"] \arrow[d] \& \Pi_{\unsure{\Lambda|\Lambda e}} \arrow[d] \\
(e,\pi') \arrow[r,"u"] \& (e,\pi)
\end{tikzcd}}
% Approximately lines 8--11.
\deleted{Hence a}
and between these diagrams there is a
transitive system of isomorphisms, hence
the diagram \((*)\) and the commutation
datum \(\alpha\).

% Approximately lines 12--19; the final note continues on scan 109.
\deleted{The link with}
When we have a \inserted{family \(\Lambda'\) (\unsure{or a set})} \deleted{\illegible{1}} of outer
homomorphisms \(\Lambda'_{i'}\ (i'\in I'):T\longrightarrow\pi'\), and a
family \(\Lambda\unsure{=}(\Lambda_i)_{i\in I}\) of outer homomorphisms \(T\longrightarrow\pi\), and some
\(\tau:I'\longrightarrow I\), \deleted{\illegible{1}}\inserted{and \((d_{i'})_{i'\in I'}\)}
\inserted{\unsure{the|some} \(d_{i'}\in\mathbb Z\), such that}
(for \(\Theta:\pi'\longrightarrow\pi\) a given outer homomorphism)
\(\forall i'\in I'\), setting \(i=\tau(i')\), we have
% The second Λ carries no prime of its own; the prime is on the index.
\[
\Lambda_i=\Lambda_{i'}\circ\Theta\circ(d_{i'}.\operatorname{id}).
\]
[\uline{N.B.} if \(I\longrightarrow\operatorname{Homext}(T,\pi)\) injective, these conditions
show that \(\tau\) is determined by \(\Theta\), and if
% Lower-left additional handwritten pagination; hand/status uncertain.
\margin{107}

\scanpage{109}
\mspage{53}
% Approximately lines 1--6; continuation from scan 108.
in addition \(\mathbb Z\) operates faithfully on \(T\), for example
if \(T\simeq\mathbb Z\), then \((d_{i'})\) is also unique \(\ldots\)]
then we construct \deleted{\unsure{B}}
% "groupoïde" stands above the struck "Produit des", "somme des" below it.
\[
\Pi_{D^*,\Lambda'}=
\deleted{\text{Product of the}}\inserted{\text{groupoid}}\ \inserted{\text{sum of the}}\ \Pi_{\Lambda'_{i'}}
\]
and the \(=\) for \(\Pi_{D^*,\Lambda}\), and we get an
\inserted{\unsure{essentially |iso}}commutative diagram of topoi
\diagram{\begin{tikzcd}[ampersand replacement=\&]
B_{D^*,\Lambda'} \arrow[r] \arrow[d] \& B_{D^*,\Lambda} \arrow[d] \\
B_{\pi'} \arrow[r] \& B_\pi
\end{tikzcd}}

% Approximately lines 7--10.
Suppose now (whew!) that
\(\Sigma\) operates on outer groups with loops \(\ldots\)
I give up --- it is obvious that
everything works fine!

% Approximately lines 11--17.
B. \uline{Return to the arithmetic case}, if we
want to describe in ``Galoisian'' terms
pairs \((X,S)\) \deleted{\unsure{connected}} anabelian \inserted{connected} over a field
\(K\) of finite type over \(\mathbb Q\). [\uline{N.B.} if we take
a \(K\) of finite type over \(\mathbb F_p\), we would have to
restrict ourselves to the fundamental groups ``prime to
\(p\)'', apart from \unsure{that|that} our developments could
be \uline{carried out} all the \(=\) \(\ldots\)]
% Lower-left additional handwritten pagination; hand/status uncertain.
\margin{108}

\scanpage{110}
% Approximately lines 1--7.
It has become clear that in terms of an
algebraic closure \(\overline K\) of \(K\), and of the
\inserted{profinite} Galois group \(\Gamma=\operatorname{Gal}(\overline K/\unsure{K|k})\), the
simplest description is in
terms of outer groups with loops, and
of outer actions of \inserted{open} subgroups \(\Sigma\subset\Gamma\)
on them. To be precise, \deleted{let} we
% Approximately lines 8--14.
choose a connected component
\(\overline X_0\) of \(\overline X\) (or what amounts to the \(=\),
\(\overline U_0\) of \(\overline U=(X\setminus S)_{\overline K}\)), let \(\Sigma\) be its
stabilizer in \(\Gamma\) (it is replaced
by a conjugate, when we change
\(\overline U_0\)). Then \(\Sigma\) operates on the scheme
\(\unsure{\overline U_0|\overline X_0}\), hence operates by outer automorphisms on its
\(\pi_1\) (considered as an outer group).
% Approximately lines 15--18.
Now on this there is a \inserted{\(T(\overline K)\)-}loop
structure, \deleted{defined by the}\inserted{\unsure{with} as} index set
\diagram{\begin{tikzcd}[ampersand replacement=\&,row sep=0.15em,column sep=0.4em]
I=S_0(\overline K)\subset \& S(\overline K) \& \qquad[\text{empty \unsure{if|iff} }S=\varnothing] \\
\& \rotatebox{90}{$\simeq$} \& \\
\& S_0(\overline K)\wedge_\Sigma\Gamma \&
\end{tikzcd}}
% Lower-left additional handwritten pagination; hand/status uncertain.
\margin{109}